\documentclass[aspectratio=169,11pt]{beamer}
\usetheme{Madrid}
\usecolortheme{default}
\usepackage[T1]{fontenc}
\setbeamertemplate{navigation symbols}{}
\setbeamertemplate{footline}[frame number]
\usepackage{enumitem}
\usepackage{amsmath,amssymb}
\usepackage{booktabs}
\usepackage{array}

\title{Databricks ML Associate}
\subtitle{High-Yield Next 100 (Set 3: Scenarios and Code Reading)}
\author{Practice Exam Deck}
\date{}

\begin{document}

\begin{frame}
\titlepage
\end{frame}

\begin{frame}{How to Use This Deck}
\begin{itemize}
\item 100 more questions (Q1--Q100), complementing Sets 1 and 2, with more scenario, calculation, and code-reading items, on topics that commonly appear on the Databricks Machine Learning Associate exam.
\item Each question slide has four options (A--D); the next slide gives the correct answer and why the others are wrong.
\item Topic order: AutoML details; code reading; Feature Engineering scenarios; Spark DataFrame code; tuning scenarios; metric calculations; ML concepts; deployment, jobs, and monitoring; Unity Catalog governance.
\end{itemize}
\end{frame}

\begin{frame}{Exam Domains Covered}
\begin{itemize}
\item \textbf{Databricks Machine Learning:} runtime, AutoML, MLflow, Unity Catalog registry, Feature Engineering client.
\item \textbf{Data Processing:} EDA, imputation, encoding, Spark ML transformers, pandas API on Spark and UDFs.
\item \textbf{Model Development:} Hyperopt, CrossValidator, metrics, bias/variance, leakage.
\item \textbf{Model Deployment:} batch, streaming, real-time serving, monitoring.
\end{itemize}
\end{frame}

% ---------------------------------------------------------------
\begin{frame}{Question 1}
Which AutoML argument lists columns that should not be used as features?
\begin{enumerate}[label=\Alph*.]
\item \texttt{skip\_cols\_list}
\item \texttt{drop\_features}
\item \texttt{exclude\_cols}
\item \texttt{ignore\_rows}
\end{enumerate}
\end{frame}
\begin{frame}{Answer 1}
\textbf{Correct: C} --- it removes columns such as IDs from training.\\
\vspace{0.5em}
\textbf{Wrong:}
\begin{itemize}
\item A: not an argument.
\item B: not an argument.
\item D: not an argument.
\end{itemize}
\end{frame}

% ---------------------------------------------------------------
\begin{frame}{Question 2}
How do you tell AutoML how long it may run?
\begin{enumerate}[label=\Alph*.]
\item Set \texttt{max\_evals}
\item Set \texttt{numFolds}
\item Set \texttt{autotermination\_minutes} on the model
\item Set \texttt{timeout\_minutes}
\end{enumerate}
\end{frame}
\begin{frame}{Answer 2}
\textbf{Correct: D} --- the experiment stops after the timeout.\\
\vspace{0.5em}
\textbf{Wrong:}
\begin{itemize}
\item A: that is a Hyperopt argument.
\item B: that is Spark ML CV.
\item C: that is a cluster setting.
\end{itemize}
\end{frame}

% ---------------------------------------------------------------
\begin{frame}{Question 3}
Which AutoML argument selects the metric used to rank trials?
\begin{enumerate}[label=\Alph*.]
\item \texttt{score\_with}
\item \texttt{primary\_metric}
\item \texttt{eval\_fn}
\item \texttt{rank\_by}
\end{enumerate}
\end{frame}
\begin{frame}{Answer 3}
\textbf{Correct: B} --- for example f1, roc\_auc, rmse, or r2.\\
\vspace{0.5em}
\textbf{Wrong:}
\begin{itemize}
\item A: not an argument.
\item C: not an argument.
\item D: not an argument.
\end{itemize}
\end{frame}

% ---------------------------------------------------------------
\begin{frame}{Question 4}
Which argument makes AutoML split classification or regression data chronologically?
\begin{enumerate}[label=\Alph*.]
\item \texttt{split\_seed}
\item \texttt{order\_by}
\item \texttt{time\_col}
\item \texttt{stratify\_col}
\end{enumerate}
\end{frame}
\begin{frame}{Answer 4}
\textbf{Correct: C} --- it uses the earliest rows for training and later rows for validation and test.\\
\vspace{0.5em}
\textbf{Wrong:}
\begin{itemize}
\item A: not an argument.
\item B: not an argument.
\item D: not an argument.
\end{itemize}
\end{frame}

% ---------------------------------------------------------------
\begin{frame}{Question 5}
Which arguments are required for \texttt{databricks.automl.forecast()} beyond the data and target?
\begin{enumerate}[label=\Alph*.]
\item \texttt{exclude\_cols} only
\item \texttt{time\_col} and \texttt{horizon} (with a frequency)
\item \texttt{pos\_label} and \texttt{threshold}
\item \texttt{numTrees} and \texttt{maxDepth}
\end{enumerate}
\end{frame}
\begin{frame}{Answer 5}
\textbf{Correct: B} --- forecasting needs the time axis and how far ahead to predict.\\
\vspace{0.5em}
\textbf{Wrong:}
\begin{itemize}
\item A: not enough for forecasting.
\item C: classification concepts.
\item D: tree hyperparameters.
\end{itemize}
\end{frame}

% ---------------------------------------------------------------
\begin{frame}{Question 6}
Which algorithm families does AutoML classification commonly try?
\begin{enumerate}[label=\Alph*.]
\item Prophet only
\item Transformers only
\item Decision trees, random forests, logistic regression, XGBoost, and LightGBM
\item k-means and PCA
\end{enumerate}
\end{frame}
\begin{frame}{Answer 6}
\textbf{Correct: C} --- AutoML tries several algorithm families.\\
\vspace{0.5em}
\textbf{Wrong:}
\begin{itemize}
\item A: a forecasting model.
\item B: not used for tabular AutoML here.
\item D: unsupervised methods.
\end{itemize}
\end{frame}

% ---------------------------------------------------------------
\begin{frame}{Question 7}
Where do AutoML trial notebooks and runs appear?
\begin{enumerate}[label=\Alph*.]
\item Only in the Spark UI
\item In a feature table
\item Only in the cluster logs
\item In an MLflow experiment
\end{enumerate}
\end{frame}
\begin{frame}{Answer 7}
\textbf{Correct: D} --- each trial is logged as a run with a notebook.\\
\vspace{0.5em}
\textbf{Wrong:}
\begin{itemize}
\item A: not stored there.
\item B: not stored there.
\item C: not stored there.
\end{itemize}
\end{frame}

% ---------------------------------------------------------------
\begin{frame}{Question 8}
A binary AutoML classification task has a rare positive class labelled "fraud". What does \texttt{pos\_label} help with?
\begin{enumerate}[label=\Alph*.]
\item Excluding the class
\item Telling AutoML which class counts as positive for precision, recall, and F1
\item Choosing the cluster size
\item Setting the number of trees
\end{enumerate}
\end{frame}
\begin{frame}{Answer 8}
\textbf{Correct: B} --- metrics depend on which class is positive.\\
\vspace{0.5em}
\textbf{Wrong:}
\begin{itemize}
\item A: not its purpose.
\item C: not its purpose.
\item D: not its purpose.
\end{itemize}
\end{frame}

% ---------------------------------------------------------------
\begin{frame}{Question 9}
What does this code do?  \texttt{with mlflow.start\_run(): mlflow.log\_param("depth", 4); mlflow.log\_metric("rmse", 2.5)}
\begin{enumerate}[label=\Alph*.]
\item Creates two separate runs
\item Registers a model
\item Creates one run holding the parameter \texttt{depth} and the metric \texttt{rmse}
\item Raises an error because params and metrics cannot be mixed
\end{enumerate}
\end{frame}
\begin{frame}{Answer 9}
\textbf{Correct: C} --- both are logged to the same active run.\\
\vspace{0.5em}
\textbf{Wrong:}
\begin{itemize}
\item A: only one run is started.
\item B: no registration call.
\item D: they can.
\end{itemize}
\end{frame}

% ---------------------------------------------------------------
\begin{frame}{Question 10}
A notebook calls \texttt{mlflow.sklearn.log\_model(model, "model")} outside any \texttt{start\_run}. What happens in a Databricks notebook?
\begin{enumerate}[label=\Alph*.]
\item It always raises an error
\item It registers to Unity Catalog
\item It logs nothing
\item A run is started automatically and the model is logged to it
\end{enumerate}
\end{frame}
\begin{frame}{Answer 10}
\textbf{Correct: D} --- MLflow auto-creates a run when logging without one.\\
\vspace{0.5em}
\textbf{Wrong:}
\begin{itemize}
\item A: an active run is created.
\item B: no registered\_model\_name was given.
\item C: it logs the model.
\end{itemize}
\end{frame}

% ---------------------------------------------------------------
\begin{frame}{Question 11}
What does \texttt{mlflow.pyfunc.load\_model("models:/main.ml.churn@Champion").predict(pdf)} return for a classifier logged with predicted labels?
\begin{enumerate}[label=\Alph*.]
\item The model signature only
\item The model's predictions for the rows in \texttt{pdf}
\item A new run ID
\item A registered model version
\end{enumerate}
\end{frame}
\begin{frame}{Answer 11}
\textbf{Correct: B} --- pyfunc models take a pandas DataFrame and return predictions.\\
\vspace{0.5em}
\textbf{Wrong:}
\begin{itemize}
\item A: not returned.
\item C: not returned.
\item D: not returned.
\end{itemize}
\end{frame}

% ---------------------------------------------------------------
\begin{frame}{Question 12}
A model was logged without a signature. What can happen when it is served?
\begin{enumerate}[label=\Alph*.]
\item It uses a GPU
\item It cannot be registered
\item It trains slower
\item Input schema validation is weaker, so mismatched columns may cause errors or silent issues
\end{enumerate}
\end{frame}
\begin{frame}{Answer 12}
\textbf{Correct: D} --- signatures help enforce expected inputs.\\
\vspace{0.5em}
\textbf{Wrong:}
\begin{itemize}
\item A: not related.
\item B: it can be registered in most cases.
\item C: not related.
\end{itemize}
\end{frame}

% ---------------------------------------------------------------
\begin{frame}{Question 13}
What is the result of \texttt{mlflow.search\_runs(order\_by=["metrics.f1 DESC"], max\_results=1)} when searching experiment runs?
\begin{enumerate}[label=\Alph*.]
\item All runs
\item The most recent run always
\item The run with the highest F1
\item The run with the lowest F1
\end{enumerate}
\end{frame}
\begin{frame}{Answer 13}
\textbf{Correct: C} --- descending order puts the highest F1 first.\\
\vspace{0.5em}
\textbf{Wrong:}
\begin{itemize}
\item A: max\_results limits to 1.
\item B: not sorted by time.
\item D: that would be ASC.
\end{itemize}
\end{frame}

% ---------------------------------------------------------------
\begin{frame}{Question 14}
Which line is correct for logging a Spark ML model?
\begin{enumerate}[label=\Alph*.]
\item \texttt{mlflow.sklearn.log\_model(pipeline\_model, "model")}
\item \texttt{mlflow.spark.save(pipeline\_model)}
\item \texttt{mlflow.spark.log\_model(pipeline\_model, "model")}
\item \texttt{mlflow.log\_dataframe(pipeline\_model)}
\end{enumerate}
\end{frame}
\begin{frame}{Answer 14}
\textbf{Correct: C} --- use the Spark flavor for Spark ML models.\\
\vspace{0.5em}
\textbf{Wrong:}
\begin{itemize}
\item A: wrong flavor.
\item B: not the logging API.
\item D: not an API.
\end{itemize}
\end{frame}

% ---------------------------------------------------------------
\begin{frame}{Question 15}
Which code builds a Spark ML pipeline correctly?
\begin{enumerate}[label=\Alph*.]
\item \texttt{Pipeline(indexer, encoder, assembler, rf)}
\item \texttt{Pipeline(stages=rf, assembler)}
\item \texttt{Pipeline.add(indexer).add(rf)}
\item \texttt{Pipeline(stages=[indexer, encoder, assembler, rf])}
\end{enumerate}
\end{frame}
\begin{frame}{Answer 15}
\textbf{Correct: D} --- stages are listed in order.\\
\vspace{0.5em}
\textbf{Wrong:}
\begin{itemize}
\item A: stages must be passed as a list.
\item B: invalid syntax.
\item C: not the API.
\end{itemize}
\end{frame}

% ---------------------------------------------------------------
\begin{frame}{Question 16}
What does this return?  \texttt{rf = RandomForestClassifier(featuresCol="features", labelCol="label"); model = rf.fit(train)}
\begin{enumerate}[label=\Alph*.]
\item A fitted \texttt{RandomForestClassificationModel}
\item A predictions DataFrame
\item A \texttt{Pipeline}
\item A metric value
\end{enumerate}
\end{frame}
\begin{frame}{Answer 16}
\textbf{Correct: A} --- fit on an estimator returns a model.\\
\vspace{0.5em}
\textbf{Wrong:}
\begin{itemize}
\item B: that is \texttt{transform}.
\item C: not returned.
\item D: not returned.
\end{itemize}
\end{frame}

% ---------------------------------------------------------------
\begin{frame}{Question 17}
What does \texttt{model.transform(test)} return for a Spark ML classifier?
\begin{enumerate}[label=\Alph*.]
\item A DataFrame with added columns such as \texttt{prediction} and \texttt{probability}
\item A list of weights
\item A metric
\item A new model
\end{enumerate}
\end{frame}
\begin{frame}{Answer 17}
\textbf{Correct: A} --- transform adds output columns.\\
\vspace{0.5em}
\textbf{Wrong:}
\begin{itemize}
\item B: not returned.
\item C: use an evaluator for that.
\item D: that is fit.
\end{itemize}
\end{frame}

% ---------------------------------------------------------------
\begin{frame}{Question 18}
Which code evaluates AUC for a Spark ML model's predictions?
\begin{enumerate}[label=\Alph*.]
\item \texttt{model.score(preds)}
\item \texttt{BinaryClassificationEvaluator(metricName="areaUnderROC").evaluate(preds)}
\item \texttt{RegressionEvaluator(metricName="areaUnderROC").evaluate(preds)}
\item \texttt{preds.auc()}
\end{enumerate}
\end{frame}
\begin{frame}{Answer 18}
\textbf{Correct: B} --- the evaluator takes the predictions DataFrame.\\
\vspace{0.5em}
\textbf{Wrong:}
\begin{itemize}
\item A: not the Spark ML API.
\item C: wrong evaluator.
\item D: not an API.
\end{itemize}
\end{frame}

% ---------------------------------------------------------------
\begin{frame}{Question 19}
Which Hyperopt call is correct for minimizing \texttt{objective} over \texttt{space} with TPE for 20 trials?
\begin{enumerate}[label=\Alph*.]
\item \texttt{tpe.fmin(space)}
\item \texttt{fmin(fn=objective, space=space, algo=tpe.suggest, max\_evals=20)}
\item \texttt{minimize(objective, space, 20)}
\item \texttt{fmin(objective, space, algo=tpe, max\_evals=20)}
\end{enumerate}
\end{frame}
\begin{frame}{Answer 19}
\textbf{Correct: B} --- fn, space, algo, and max\_evals are the core arguments.\\
\vspace{0.5em}
\textbf{Wrong:}
\begin{itemize}
\item A: not the API.
\item C: not the API.
\item D: algo needs \texttt{tpe.suggest}.
\end{itemize}
\end{frame}

% ---------------------------------------------------------------
\begin{frame}{Question 20}
An objective returns \texttt{-accuracy} while the team reports the best result. What should they report?
\begin{enumerate}[label=\Alph*.]
\item The number of trials
\item The loss as the accuracy directly
\item The worst loss
\item The negated best loss, which is the best accuracy
\end{enumerate}
\end{frame}
\begin{frame}{Answer 20}
\textbf{Correct: D} --- loss is -accuracy, so flip the sign.\\
\vspace{0.5em}
\textbf{Wrong:}
\begin{itemize}
\item A: not a metric.
\item B: the sign is flipped.
\item C: not the best.
\end{itemize}
\end{frame}

% ---------------------------------------------------------------
\begin{frame}{Question 21}
A scikit-learn Hyperopt objective is run with \texttt{SparkTrials(parallelism=8)} on an 8-worker cluster. What is the main benefit?
\begin{enumerate}[label=\Alph*.]
\item Up to 8 trials can run at the same time on workers
\item TPE becomes random search
\item Each model trains on all workers
\item Trials are guaranteed to be better
\end{enumerate}
\end{frame}
\begin{frame}{Answer 21}
\textbf{Correct: A} --- SparkTrials distributes single-node trials.\\
\vspace{0.5em}
\textbf{Wrong:}
\begin{itemize}
\item B: TPE still runs.
\item C: sklearn does not distribute.
\item D: not guaranteed.
\end{itemize}
\end{frame}

% ---------------------------------------------------------------
\begin{frame}{Question 22}
With SparkTrials, how are MLflow runs for trials typically organized?
\begin{enumerate}[label=\Alph*.]
\item One run per cluster node
\item A parent run with a child run per trial
\item One run per Spark task
\item No runs are created
\end{enumerate}
\end{frame}
\begin{frame}{Answer 22}
\textbf{Correct: B} --- Hyperopt on Databricks logs trials as nested runs.\\
\vspace{0.5em}
\textbf{Wrong:}
\begin{itemize}
\item A: not how it is organized.
\item C: not how it is organized.
\item D: runs are created.
\end{itemize}
\end{frame}

% ---------------------------------------------------------------
\begin{frame}{Question 23}
Which snippet pulls the Python value of the best run's RMSE from a \texttt{search\_runs} DataFrame \texttt{runs}?
\begin{enumerate}[label=\Alph*.]
\item \texttt{runs.get\_metric()}
\item \texttt{runs.rmse()}
\item \texttt{runs["rmse"].best}
\item \texttt{runs.loc[0, "metrics.rmse"]}
\end{enumerate}
\end{frame}
\begin{frame}{Answer 23}
\textbf{Correct: D} --- metrics columns are named \texttt{metrics.<name>}.\\
\vspace{0.5em}
\textbf{Wrong:}
\begin{itemize}
\item A: not valid.
\item B: not valid.
\item C: not valid.
\end{itemize}
\end{frame}

% ---------------------------------------------------------------
\begin{frame}{Question 24}
What does \texttt{df.randomSplit([0.7, 0.3], seed=1)} return?
\begin{enumerate}[label=\Alph*.]
\item A dict of metrics
\item A list of two DataFrames
\item A Pipeline
\item A single DataFrame
\end{enumerate}
\end{frame}
\begin{frame}{Answer 24}
\textbf{Correct: B} --- the weights define the split.\\
\vspace{0.5em}
\textbf{Wrong:}
\begin{itemize}
\item A: not returned.
\item C: not returned.
\item D: it returns multiple.
\end{itemize}
\end{frame}

% ---------------------------------------------------------------
\begin{frame}{Question 25}
Which code removes the \texttt{id} column and keeps the rest?
\begin{enumerate}[label=\Alph*.]
\item \texttt{df.select("id")}
\item \texttt{df.drop("id")}
\item \texttt{df.remove("id")}
\item \texttt{df.filter("id")}
\end{enumerate}
\end{frame}
\begin{frame}{Answer 25}
\textbf{Correct: B} --- drop removes the named column.\\
\vspace{0.5em}
\textbf{Wrong:}
\begin{itemize}
\item A: keeps only id.
\item C: not an API.
\item D: filters rows.
\end{itemize}
\end{frame}

% ---------------------------------------------------------------
\begin{frame}{Question 26}
Which code keeps rows where \texttt{age} is greater than 30?
\begin{enumerate}[label=\Alph*.]
\item \texttt{df.limit(df.age > 30)}
\item \texttt{df.select(df.age > 30)}
\item \texttt{df.drop(df.age > 30)}
\item \texttt{df.filter(df.age > 30)}
\end{enumerate}
\end{frame}
\begin{frame}{Answer 26}
\textbf{Correct: D} --- filter keeps rows matching the condition.\\
\vspace{0.5em}
\textbf{Wrong:}
\begin{itemize}
\item A: limit takes a number.
\item B: returns a boolean column.
\item C: drops columns.
\end{itemize}
\end{frame}

% ---------------------------------------------------------------
\begin{frame}{Question 27}
Which code casts \texttt{price} to double?
\begin{enumerate}[label=\Alph*.]
\item \texttt{df.withColumn("price", df.price.cast("double"))}
\item \texttt{df.price.toDouble()}
\item \texttt{df.drop("price", "double")}
\item \texttt{df.convert("price", "double")}
\end{enumerate}
\end{frame}
\begin{frame}{Answer 27}
\textbf{Correct: A} --- cast changes the column type.\\
\vspace{0.5em}
\textbf{Wrong:}
\begin{itemize}
\item B: not an API.
\item C: drops columns.
\item D: not an API.
\end{itemize}
\end{frame}

% ---------------------------------------------------------------
\begin{frame}{Question 28}
Which code creates a Spark DataFrame from a pandas DataFrame \texttt{pdf}?
\begin{enumerate}[label=\Alph*.]
\item \texttt{spark.createDataFrame(pdf)}
\item \texttt{pdf.toSpark()}
\item \texttt{spark.load(pdf)}
\item \texttt{pdf.to\_delta()}
\end{enumerate}
\end{frame}
\begin{frame}{Answer 28}
\textbf{Correct: A} --- it converts pandas to Spark.\\
\vspace{0.5em}
\textbf{Wrong:}
\begin{itemize}
\item B: not an API.
\item C: not an API.
\item D: not the conversion.
\end{itemize}
\end{frame}

% ---------------------------------------------------------------
\begin{frame}{Question 29}
Training features were computed with a window using future data relative to each label date. What is the problem?
\begin{enumerate}[label=\Alph*.]
\item Slow inference
\item Underfitting only
\item Too few trees
\item Data leakage that makes offline metrics too optimistic
\end{enumerate}
\end{frame}
\begin{frame}{Answer 29}
\textbf{Correct: D} --- the model sees information unavailable at prediction time.\\
\vspace{0.5em}
\textbf{Wrong:}
\begin{itemize}
\item A: not the problem.
\item B: not the main problem.
\item C: not the problem.
\end{itemize}
\end{frame}

% ---------------------------------------------------------------
\begin{frame}{Question 30}
A feature table has \texttt{customer\_id} as key and a \texttt{ts} timestamp key. The label table has \texttt{customer\_id} and \texttt{event\_time}. What is needed for a point-in-time join?
\begin{enumerate}[label=\Alph*.]
\item \texttt{lookup\_key="ts"} only
\item A \texttt{FeatureLookup} with \texttt{timestamp\_lookup\_key="event\_time"}
\item A cross join
\item Rounding all dates
\end{enumerate}
\end{frame}
\begin{frame}{Answer 30}
\textbf{Correct: B} --- this joins the latest feature values at or before each event time.\\
\vspace{0.5em}
\textbf{Wrong:}
\begin{itemize}
\item A: ignores the label time column.
\item C: causes leakage and row explosion.
\item D: does not fix joins.
\end{itemize}
\end{frame}

% ---------------------------------------------------------------
\begin{frame}{Question 31}
A team wants to add a new column to an existing feature table without overwriting existing rows. What should they do?
\begin{enumerate}[label=\Alph*.]
\item Create a new cluster
\item Drop and recreate the table every time
\item Use \texttt{mode="overwrite"} with only new rows
\item Write the new column with \texttt{fe.write\_table(..., mode="merge")} keyed on the primary key
\end{enumerate}
\end{frame}
\begin{frame}{Answer 31}
\textbf{Correct: D} --- merge adds or updates by key.\\
\vspace{0.5em}
\textbf{Wrong:}
\begin{itemize}
\item A: not related.
\item B: loses history.
\item C: replaces data.
\end{itemize}
\end{frame}

% ---------------------------------------------------------------
\begin{frame}{Question 32}
Training uses \texttt{FeatureLookup}s over two feature tables. Which column must exist in the label DataFrame?
\begin{enumerate}[label=\Alph*.]
\item The lookup key columns (and a timestamp column for point-in-time lookups)
\item A copy of every feature
\item Only the label column
\item An MLflow run ID
\end{enumerate}
\end{frame}
\begin{frame}{Answer 32}
\textbf{Correct: A} --- keys connect the label rows to feature rows.\\
\vspace{0.5em}
\textbf{Wrong:}
\begin{itemize}
\item B: unnecessary.
\item C: keys are required too.
\item D: not needed.
\end{itemize}
\end{frame}

% ---------------------------------------------------------------
\begin{frame}{Question 33}
Why might offline model accuracy be good but online accuracy poor?
\begin{enumerate}[label=\Alph*.]
\item The cluster is too small
\item The model has too few parameters
\item Features are computed differently online than offline (training-serving skew)
\item MLflow is not installed
\end{enumerate}
\end{frame}
\begin{frame}{Answer 33}
\textbf{Correct: C} --- a feature store reduces this risk.\\
\vspace{0.5em}
\textbf{Wrong:}
\begin{itemize}
\item A: not implied.
\item B: not implied.
\item D: not implied.
\end{itemize}
\end{frame}

% ---------------------------------------------------------------
\begin{frame}{Question 34}
Which key type does \texttt{fe.create\_table} need so that rows can be uniquely identified?
\begin{enumerate}[label=\Alph*.]
\item \texttt{foreign\_keys}
\item \texttt{cluster\_keys} only
\item \texttt{primary\_keys}
\item \texttt{partition\_keys} only
\end{enumerate}
\end{frame}
\begin{frame}{Answer 34}
\textbf{Correct: C} --- primary keys identify feature rows.\\
\vspace{0.5em}
\textbf{Wrong:}
\begin{itemize}
\item A: not required.
\item B: for layout only.
\item D: for layout only.
\end{itemize}
\end{frame}

% ---------------------------------------------------------------
\begin{frame}{Question 35}
A feature table is large and only some features are needed for a model. How do you limit the features used?
\begin{enumerate}[label=\Alph*.]
\item Drop the table
\item Use \texttt{df.limit(1)}
\item Specify \texttt{feature\_names} in the \texttt{FeatureLookup}
\item Disable lineage
\end{enumerate}
\end{frame}
\begin{frame}{Answer 35}
\textbf{Correct: C} --- only the listed features are joined.\\
\vspace{0.5em}
\textbf{Wrong:}
\begin{itemize}
\item A: loses features.
\item B: reduces rows, not features.
\item D: unrelated.
\end{itemize}
\end{frame}

% ---------------------------------------------------------------
\begin{frame}{Question 36}
What is true about a model logged by \texttt{fe.log\_model()} when scored with \texttt{fe.score\_batch()}?
\begin{enumerate}[label=\Alph*.]
\item All features must be in the input DataFrame
\item The model is retrained
\item Feature values are looked up from the feature tables by key
\item Scoring always uses the training data
\end{enumerate}
\end{frame}
\begin{frame}{Answer 36}
\textbf{Correct: C} --- the input DataFrame only needs keys and non-lookup inputs.\\
\vspace{0.5em}
\textbf{Wrong:}
\begin{itemize}
\item A: that defeats the lookup.
\item B: not true.
\item D: not true.
\end{itemize}
\end{frame}

% ---------------------------------------------------------------
\begin{frame}{Question 37}
Which code adds a column \texttt{is\_high} that is 1 when \texttt{amount > 100} and 0 otherwise?
\begin{enumerate}[label=\Alph*.]
\item \texttt{df.withColumn("is\_high", F.when(df.amount > 100, 1).otherwise(0))}
\item \texttt{df.withColumn("is\_high").when(1)}
\item \texttt{df.add("is\_high", df.amount > 100)}
\item \texttt{df.if("amount", 100)}
\end{enumerate}
\end{frame}
\begin{frame}{Answer 37}
\textbf{Correct: A} --- when/otherwise builds conditional columns.\\
\vspace{0.5em}
\textbf{Wrong:}
\begin{itemize}
\item B: invalid.
\item C: not an API.
\item D: not an API.
\end{itemize}
\end{frame}

% ---------------------------------------------------------------
\begin{frame}{Question 38}
Which aggregation computes the average \texttt{price} per \texttt{region}?
\begin{enumerate}[label=\Alph*.]
\item \texttt{df.avg("region", "price")}
\item \texttt{df.agg("region").avg("price")}
\item \texttt{df.groupBy("price").agg(F.avg("region"))}
\item \texttt{df.groupBy("region").agg(F.avg("price"))}
\end{enumerate}
\end{frame}
\begin{frame}{Answer 38}
\textbf{Correct: D} --- groupBy then agg.\\
\vspace{0.5em}
\textbf{Wrong:}
\begin{itemize}
\item A: not an API.
\item B: invalid.
\item C: swapped columns.
\end{itemize}
\end{frame}

% ---------------------------------------------------------------
\begin{frame}{Question 39}
Which join keeps all rows from the left DataFrame, with nulls where the right has no match?
\begin{enumerate}[label=\Alph*.]
\item Left join
\item Inner join
\item Cross join
\item Left anti join
\end{enumerate}
\end{frame}
\begin{frame}{Answer 39}
\textbf{Correct: A} --- left join preserves all left rows.\\
\vspace{0.5em}
\textbf{Wrong:}
\begin{itemize}
\item B: keeps only matches.
\item C: returns all pairs.
\item D: keeps only non-matching rows.
\end{itemize}
\end{frame}

% ---------------------------------------------------------------
\begin{frame}{Question 40}
Which join finds left rows with no match in the right DataFrame?
\begin{enumerate}[label=\Alph*.]
\item Left anti join
\item Inner join
\item Full outer join
\item Left semi join
\end{enumerate}
\end{frame}
\begin{frame}{Answer 40}
\textbf{Correct: A} --- it returns non-matching left rows.\\
\vspace{0.5em}
\textbf{Wrong:}
\begin{itemize}
\item B: returns matches.
\item C: returns all rows.
\item D: returns matching left rows.
\end{itemize}
\end{frame}

% ---------------------------------------------------------------
\begin{frame}{Question 41}
Which call stacks two DataFrames with the same schema vertically?
\begin{enumerate}[label=\Alph*.]
\item \texttt{df1.union(df2)}
\item \texttt{df1.stack(df2)}
\item \texttt{df1.join(df2)}
\item \texttt{df1.merge(df2)}
\end{enumerate}
\end{frame}
\begin{frame}{Answer 41}
\textbf{Correct: A} --- union appends rows.\\
\vspace{0.5em}
\textbf{Wrong:}
\begin{itemize}
\item B: not an API.
\item C: combines columns by key.
\item D: not a Spark DataFrame method.
\end{itemize}
\end{frame}

% ---------------------------------------------------------------
\begin{frame}{Question 42}
Which code orders a Spark DataFrame by \texttt{score} highest first?
\begin{enumerate}[label=\Alph*.]
\item \texttt{df.limit("score")}
\item \texttt{df.sortUp("score")}
\item \texttt{df.orderBy(F.col("score").desc())}
\item \texttt{df.orderBy("score", asc=True)}
\end{enumerate}
\end{frame}
\begin{frame}{Answer 42}
\textbf{Correct: C} --- desc sorts in descending order.\\
\vspace{0.5em}
\textbf{Wrong:}
\begin{itemize}
\item A: limit takes a count.
\item B: not an API.
\item D: ascending.
\end{itemize}
\end{frame}

% ---------------------------------------------------------------
\begin{frame}{Question 43}
A pandas-on-Spark DataFrame \texttt{psdf} is needed as a Spark DataFrame again. Which call converts it?
\begin{enumerate}[label=\Alph*.]
\item \texttt{spark.from\_pandas(psdf)}
\item \texttt{psdf.collect\_spark()}
\item \texttt{psdf.to\_spark()}
\item \texttt{psdf.toDF()}
\end{enumerate}
\end{frame}
\begin{frame}{Answer 43}
\textbf{Correct: C} --- to\_spark converts to a Spark DataFrame.\\
\vspace{0.5em}
\textbf{Wrong:}
\begin{itemize}
\item A: not an API.
\item B: not an API.
\item D: not the pandas API on Spark method.
\end{itemize}
\end{frame}

% ---------------------------------------------------------------
\begin{frame}{Question 44}
A scikit-learn model must score a very large Spark DataFrame. Which approach distributes scoring?
\begin{enumerate}[label=\Alph*.]
\item \texttt{df.coalesce(1)} and loop
\item \texttt{df.cache()} only
\item A pandas UDF or \texttt{mlflow.pyfunc.spark\_udf} wrapping the model
\item \texttt{df.toPandas().predict()}
\end{enumerate}
\end{frame}
\begin{frame}{Answer 44}
\textbf{Correct: C} --- both run predictions on workers.\\
\vspace{0.5em}
\textbf{Wrong:}
\begin{itemize}
\item A: not distributed.
\item B: does not score.
\item D: runs on the driver only.
\end{itemize}
\end{frame}

% ---------------------------------------------------------------
\begin{frame}{Question 45}
Which Spark SQL statement creates a table from a query result?
\begin{enumerate}[label=\Alph*.]
\item \texttt{CREATE TABLE t AS SELECT ...}
\item \texttt{MAKE TABLE t FROM ...}
\item \texttt{COPY TABLE t}
\item \texttt{SELECT INTO t}
\end{enumerate}
\end{frame}
\begin{frame}{Answer 45}
\textbf{Correct: A} --- CTAS creates a table from a query.\\
\vspace{0.5em}
\textbf{Wrong:}
\begin{itemize}
\item B: not valid.
\item C: not valid.
\item D: not Spark SQL.
\end{itemize}
\end{frame}

% ---------------------------------------------------------------
\begin{frame}{Question 46}
What does \texttt{display(df)} offer in a Databricks notebook?
\begin{enumerate}[label=\Alph*.]
\item A deployed endpoint
\item A trained model
\item A Git commit
\item An interactive table with built-in charting options
\end{enumerate}
\end{frame}
\begin{frame}{Answer 46}
\textbf{Correct: D} --- it renders data with plots.\\
\vspace{0.5em}
\textbf{Wrong:}
\begin{itemize}
\item A: not provided.
\item B: not provided.
\item C: not provided.
\end{itemize}
\end{frame}

% ---------------------------------------------------------------
\begin{frame}{Question 47}
A 5-fold CrossValidator uses a grid of 2 x 3 x 2 combinations. How many models are trained during the search (excluding the final refit)?
\begin{enumerate}[label=\Alph*.]
\item 60
\item 12
\item 30
\item 17
\end{enumerate}
\end{frame}
\begin{frame}{Answer 47}
\textbf{Correct: A} --- 12 combinations times 5 folds.\\
\vspace{0.5em}
\textbf{Wrong:}
\begin{itemize}
\item B: ignores the folds.
\item C: miscounts the grid or folds.
\item D: adds instead of multiplying.
\end{itemize}
\end{frame}

% ---------------------------------------------------------------
\begin{frame}{Question 48}
A team has limited compute and many hyperparameters. Which strategy is usually more efficient than full grid search?
\begin{enumerate}[label=\Alph*.]
\item Removing validation
\item Random search or Bayesian optimization (Hyperopt/TPE)
\item Training on the test set
\item Exhaustive grid with finer steps
\end{enumerate}
\end{frame}
\begin{frame}{Answer 48}
\textbf{Correct: B} --- they explore the space with fewer trials.\\
\vspace{0.5em}
\textbf{Wrong:}
\begin{itemize}
\item A: invalid.
\item C: invalid.
\item D: more costly.
\end{itemize}
\end{frame}

% ---------------------------------------------------------------
\begin{frame}{Question 49}
A tuning job finds the best validation score at the boundary of the search range. What should the team do?
\begin{enumerate}[label=\Alph*.]
\item Expand the search range in that direction and rerun
\item Accept it as globally optimal
\item Shrink the range
\item Switch to the test set
\end{enumerate}
\end{frame}
\begin{frame}{Answer 49}
\textbf{Correct: A} --- the optimum may lie outside the range.\\
\vspace{0.5em}
\textbf{Wrong:}
\begin{itemize}
\item B: may be wrong.
\item C: makes it worse.
\item D: invalid.
\end{itemize}
\end{frame}

% ---------------------------------------------------------------
\begin{frame}{Question 50}
Why is the test set not used for hyperparameter tuning?
\begin{enumerate}[label=\Alph*.]
\item Using it would bias the final performance estimate
\item MLflow forbids it
\item It is too small to compute metrics
\item Hyperopt cannot read it
\end{enumerate}
\end{frame}
\begin{frame}{Answer 50}
\textbf{Correct: A} --- tuning on test data leaks information.\\
\vspace{0.5em}
\textbf{Wrong:}
\begin{itemize}
\item B: not the reason.
\item C: not the reason.
\item D: not the reason.
\end{itemize}
\end{frame}

% ---------------------------------------------------------------
\begin{frame}{Question 51}
What does \texttt{TrainValidationSplit} need compared with \texttt{CrossValidator}?
\begin{enumerate}[label=\Alph*.]
\item A Delta trigger
\item A serving endpoint
\item A feature table
\item A \texttt{trainRatio} instead of \texttt{numFolds}
\end{enumerate}
\end{frame}
\begin{frame}{Answer 51}
\textbf{Correct: D} --- trainRatio sets the split proportion.\\
\vspace{0.5em}
\textbf{Wrong:}
\begin{itemize}
\item A: not required.
\item B: not required.
\item C: not required.
\end{itemize}
\end{frame}

% ---------------------------------------------------------------
\begin{frame}{Question 52}
A tuned model performs well in CV but poorly on the held-out test set. What is a likely cause?
\begin{enumerate}[label=\Alph*.]
\item The test set is perfect
\item MLflow logged too many metrics
\item Too many workers
\item Overfitting to the validation process or leakage during CV
\end{enumerate}
\end{frame}
\begin{frame}{Answer 52}
\textbf{Correct: D} --- tuning can overfit validation folds.\\
\vspace{0.5em}
\textbf{Wrong:}
\begin{itemize}
\item A: not implied.
\item B: not related.
\item C: not related.
\end{itemize}
\end{frame}

% ---------------------------------------------------------------
\begin{frame}{Question 53}
Where should a Spark ML \texttt{Pipeline} containing preprocessing be placed when tuning with \texttt{CrossValidator}?
\begin{enumerate}[label=\Alph*.]
\item Run only on the test data
\item As the estimator inside the CrossValidator so preprocessing is refit per fold
\item Skip preprocessing
\item Fit once on all data before CV
\end{enumerate}
\end{frame}
\begin{frame}{Answer 53}
\textbf{Correct: B} --- this avoids leakage across folds.\\
\vspace{0.5em}
\textbf{Wrong:}
\begin{itemize}
\item A: invalid.
\item C: not required.
\item D: leaks fold information.
\end{itemize}
\end{frame}

% ---------------------------------------------------------------
\begin{frame}{Question 54}
Hyperopt trials all return the same loss. What is a likely cause?
\begin{enumerate}[label=\Alph*.]
\item The cluster is too large
\item The dataset is too small
\item TPE is too random
\item The objective ignores the sampled hyperparameters
\end{enumerate}
\end{frame}
\begin{frame}{Answer 54}
\textbf{Correct: D} --- the parameters are not passed into the model.\\
\vspace{0.5em}
\textbf{Wrong:}
\begin{itemize}
\item A: not related.
\item B: not necessarily related.
\item C: does not explain identical loss.
\end{itemize}
\end{frame}

% ---------------------------------------------------------------
\begin{frame}{Question 55}
A classifier has TP=40, FP=10, FN=20, TN=30. What is its precision?
\begin{enumerate}[label=\Alph*.]
\item 0.67
\item 0.70
\item 0.50
\item 0.80
\end{enumerate}
\end{frame}
\begin{frame}{Answer 55}
\textbf{Correct: D} --- 40 / (40 + 10).\\
\vspace{0.5em}
\textbf{Wrong:}
\begin{itemize}
\item A: that is recall, 40/60.
\item B: that is accuracy, 70/100.
\item C: incorrect.
\end{itemize}
\end{frame}

% ---------------------------------------------------------------
\begin{frame}{Question 56}
A classifier has TP=40, FP=10, FN=20, TN=30. What is its recall?
\begin{enumerate}[label=\Alph*.]
\item 0.50
\item 0.30
\item 0.80
\item About 0.67
\end{enumerate}
\end{frame}
\begin{frame}{Answer 56}
\textbf{Correct: D} --- 40 / (40 + 20).\\
\vspace{0.5em}
\textbf{Wrong:}
\begin{itemize}
\item A: incorrect.
\item B: incorrect.
\item C: that is precision.
\end{itemize}
\end{frame}

% ---------------------------------------------------------------
\begin{frame}{Question 57}
A classifier has TP=40, FP=10, FN=20, TN=30. What is its accuracy?
\begin{enumerate}[label=\Alph*.]
\item 0.67
\item 0.80
\item 0.40
\item 0.70
\end{enumerate}
\end{frame}
\begin{frame}{Answer 57}
\textbf{Correct: D} --- (40 + 30) / 100.\\
\vspace{0.5em}
\textbf{Wrong:}
\begin{itemize}
\item A: that is recall.
\item B: that is precision.
\item C: incorrect.
\end{itemize}
\end{frame}

% ---------------------------------------------------------------
\begin{frame}{Question 58}
Prediction errors are 1, 1, 1, and 5. What is the MAE?
\begin{enumerate}[label=\Alph*.]
\item 2.0
\item 8.0
\item 2.65
\item 1.0
\end{enumerate}
\end{frame}
\begin{frame}{Answer 58}
\textbf{Correct: A} --- (1 + 1 + 1 + 5) / 4.\\
\vspace{0.5em}
\textbf{Wrong:}
\begin{itemize}
\item B: that is the sum.
\item C: that is the RMSE.
\item D: ignores the large error.
\end{itemize}
\end{frame}

% ---------------------------------------------------------------
\begin{frame}{Question 59}
Prediction errors are 1, 1, 1, and 5. Which is larger, RMSE or MAE, and why?
\begin{enumerate}[label=\Alph*.]
\item RMSE, because squaring weights the large error more
\item MAE, because it sums errors
\item Neither can be computed
\item They are always equal
\end{enumerate}
\end{frame}
\begin{frame}{Answer 59}
\textbf{Correct: A} --- RMSE is about 2.65 versus MAE 2.0.\\
\vspace{0.5em}
\textbf{Wrong:}
\begin{itemize}
\item B: MAE averages absolute errors.
\item C: both can be computed.
\item D: they differ.
\end{itemize}
\end{frame}

% ---------------------------------------------------------------
\begin{frame}{Question 60}
A bank wants to minimize wrongly flagging good customers as fraud. Which metric focuses on this?
\begin{enumerate}[label=\Alph*.]
\item Recall
\item R2
\item Precision
\item Silhouette
\end{enumerate}
\end{frame}
\begin{frame}{Answer 60}
\textbf{Correct: C} --- precision penalizes false positives.\\
\vspace{0.5em}
\textbf{Wrong:}
\begin{itemize}
\item A: focuses on false negatives.
\item B: regression metric.
\item D: clustering metric.
\end{itemize}
\end{frame}

% ---------------------------------------------------------------
\begin{frame}{Question 61}
A ranking-quality metric that is threshold independent for binary classification is:
\begin{enumerate}[label=\Alph*.]
\item ROC AUC
\item Precision at 0.5
\item Accuracy at 0.5
\item F1 at 0.5
\end{enumerate}
\end{frame}
\begin{frame}{Answer 61}
\textbf{Correct: A} --- it summarizes performance over all thresholds.\\
\vspace{0.5em}
\textbf{Wrong:}
\begin{itemize}
\item B: threshold dependent.
\item C: threshold dependent.
\item D: threshold dependent.
\end{itemize}
\end{frame}

% ---------------------------------------------------------------
\begin{frame}{Question 62}
For highly imbalanced data, which curve is often more informative than the ROC curve?
\begin{enumerate}[label=\Alph*.]
\item The residual plot
\item The precision-recall curve
\item The learning-rate curve
\item The elbow curve
\end{enumerate}
\end{frame}
\begin{frame}{Answer 62}
\textbf{Correct: B} --- it focuses on the positive class.\\
\vspace{0.5em}
\textbf{Wrong:}
\begin{itemize}
\item A: for regression.
\item C: not a standard evaluation.
\item D: for clustering.
\end{itemize}
\end{frame}

% ---------------------------------------------------------------
\begin{frame}{Question 63}
When averaging F1 across classes in multiclass classification, what does macro averaging do?
\begin{enumerate}[label=\Alph*.]
\item Pools all predictions first
\item Ignores rare classes
\item Computes F1 per class and averages them equally
\item Weights classes by frequency
\end{enumerate}
\end{frame}
\begin{frame}{Answer 63}
\textbf{Correct: C} --- each class counts the same.\\
\vspace{0.5em}
\textbf{Wrong:}
\begin{itemize}
\item A: that is micro averaging.
\item B: not what it does.
\item D: that is weighted averaging.
\end{itemize}
\end{frame}

% ---------------------------------------------------------------
\begin{frame}{Question 64}
A model with R2 = -0.2 on the test set indicates:
\begin{enumerate}[label=\Alph*.]
\item No variance in the labels
\item It is worse than predicting the mean
\item A perfect fit
\item A bug in the evaluator only
\end{enumerate}
\end{frame}
\begin{frame}{Answer 64}
\textbf{Correct: B} --- R2 can be negative on held-out data.\\
\vspace{0.5em}
\textbf{Wrong:}
\begin{itemize}
\item A: not implied.
\item C: that is near 1.
\item D: not necessarily.
\end{itemize}
\end{frame}

% ---------------------------------------------------------------
\begin{frame}{Question 65}
Which is an example of supervised learning?
\begin{enumerate}[label=\Alph*.]
\item Reducing dimensions with PCA
\item Grouping customers without labels
\item Detecting topics without labels
\item Predicting house prices from labelled sales data
\end{enumerate}
\end{frame}
\begin{frame}{Answer 65}
\textbf{Correct: D} --- labels guide the training.\\
\vspace{0.5em}
\textbf{Wrong:}
\begin{itemize}
\item A: unsupervised.
\item B: unsupervised.
\item C: unsupervised.
\end{itemize}
\end{frame}

% ---------------------------------------------------------------
\begin{frame}{Question 66}
Which is an example of unsupervised learning?
\begin{enumerate}[label=\Alph*.]
\item Classifying spam from labelled email
\item Clustering customers into segments without labels
\item Forecasting sales from history
\item Predicting churn from labels
\end{enumerate}
\end{frame}
\begin{frame}{Answer 66}
\textbf{Correct: B} --- no target is used.\\
\vspace{0.5em}
\textbf{Wrong:}
\begin{itemize}
\item A: supervised.
\item C: supervised.
\item D: supervised.
\end{itemize}
\end{frame}

% ---------------------------------------------------------------
\begin{frame}{Question 67}
Which algorithm classifies a point by the majority label of its nearest neighbors?
\begin{enumerate}[label=\Alph*.]
\item k-means
\item k-nearest neighbors
\item PCA
\item Linear regression
\end{enumerate}
\end{frame}
\begin{frame}{Answer 67}
\textbf{Correct: B} --- it uses distance to neighbors.\\
\vspace{0.5em}
\textbf{Wrong:}
\begin{itemize}
\item A: unsupervised clustering.
\item C: dimensionality reduction.
\item D: predicts continuous values.
\end{itemize}
\end{frame}

% ---------------------------------------------------------------
\begin{frame}{Question 68}
Why does k-nearest neighbors usually need feature scaling?
\begin{enumerate}[label=\Alph*.]
\item Distances are dominated by features with large numeric ranges
\item It does not use distances
\item It uses only categorical features
\item It cannot handle numbers
\end{enumerate}
\end{frame}
\begin{frame}{Answer 68}
\textbf{Correct: A} --- scaling puts features on comparable scales.\\
\vspace{0.5em}
\textbf{Wrong:}
\begin{itemize}
\item B: it does.
\item C: not true.
\item D: it can.
\end{itemize}
\end{frame}

% ---------------------------------------------------------------
\begin{frame}{Question 69}
What is a main advantage of a random forest over a single decision tree?
\begin{enumerate}[label=\Alph*.]
\item It needs no data
\item Lower variance and less overfitting through averaging many trees
\item It never overfits
\item It is fully interpretable
\end{enumerate}
\end{frame}
\begin{frame}{Answer 69}
\textbf{Correct: B} --- bagging reduces variance.\\
\vspace{0.5em}
\textbf{Wrong:}
\begin{itemize}
\item A: not true.
\item C: it can still overfit.
\item D: a single tree is more interpretable.
\end{itemize}
\end{frame}

% ---------------------------------------------------------------
\begin{frame}{Question 70}
In gradient boosting, how are trees combined?
\begin{enumerate}[label=\Alph*.]
\item In parallel as independent trees only
\item Sequentially, each tree fitting the errors of the previous ensemble
\item By clustering features
\item By averaging identical trees
\end{enumerate}
\end{frame}
\begin{frame}{Answer 70}
\textbf{Correct: B} --- boosting corrects residual errors step by step.\\
\vspace{0.5em}
\textbf{Wrong:}
\begin{itemize}
\item A: that is bagging.
\item C: not boosting.
\item D: not boosting.
\end{itemize}
\end{frame}

% ---------------------------------------------------------------
\begin{frame}{Question 71}
Which scenario describes underfitting?
\begin{enumerate}[label=\Alph*.]
\item Perfect test accuracy
\item Training stops early by design
\item Low accuracy on both training and validation data
\item High training accuracy, low validation accuracy
\end{enumerate}
\end{frame}
\begin{frame}{Answer 71}
\textbf{Correct: C} --- the model is too simple to learn the pattern.\\
\vspace{0.5em}
\textbf{Wrong:}
\begin{itemize}
\item A: not underfitting.
\item B: not underfitting.
\item D: that is overfitting.
\end{itemize}
\end{frame}

% ---------------------------------------------------------------
\begin{frame}{Question 72}
Which action most directly reduces overfitting?
\begin{enumerate}[label=\Alph*.]
\item Train on the test set
\item Remove the validation set
\item Add regularization or use more training data
\item Increase model complexity
\end{enumerate}
\end{frame}
\begin{frame}{Answer 72}
\textbf{Correct: C} --- both reduce variance.\\
\vspace{0.5em}
\textbf{Wrong:}
\begin{itemize}
\item A: invalid.
\item B: removes the signal.
\item D: can increase overfitting.
\end{itemize}
\end{frame}

% ---------------------------------------------------------------
\begin{frame}{Question 73}
Which action can reduce underfitting?
\begin{enumerate}[label=\Alph*.]
\item Add stronger regularization
\item Use a more flexible model or add informative features
\item Shuffle labels
\item Reduce training data
\end{enumerate}
\end{frame}
\begin{frame}{Answer 73}
\textbf{Correct: B} --- it increases the model's capacity.\\
\vspace{0.5em}
\textbf{Wrong:}
\begin{itemize}
\item A: can worsen underfitting.
\item C: destroys signal.
\item D: does not help.
\end{itemize}
\end{frame}

% ---------------------------------------------------------------
\begin{frame}{Question 74}
Why is stratified k-fold useful for classification?
\begin{enumerate}[label=\Alph*.]
\item It speeds up training by ignoring labels
\item It removes the minority class
\item It shuffles labels
\item Each fold keeps roughly the same class proportions
\end{enumerate}
\end{frame}
\begin{frame}{Answer 74}
\textbf{Correct: D} --- this stabilizes evaluation on imbalanced data.\\
\vspace{0.5em}
\textbf{Wrong:}
\begin{itemize}
\item A: it does not.
\item B: it does not.
\item C: it does not.
\end{itemize}
\end{frame}

% ---------------------------------------------------------------
\begin{frame}{Question 75}
A feature is the label's exact copy under a different name. What is this called?
\begin{enumerate}[label=\Alph*.]
\item Cross-validation
\item Regularization
\item Target leakage
\item Normalization
\end{enumerate}
\end{frame}
\begin{frame}{Answer 75}
\textbf{Correct: C} --- the feature reveals the answer.\\
\vspace{0.5em}
\textbf{Wrong:}
\begin{itemize}
\item A: not leakage.
\item B: not leakage.
\item D: not leakage.
\end{itemize}
\end{frame}

% ---------------------------------------------------------------
\begin{frame}{Question 76}
Logistic regression outputs probabilities. How do you control the trade-off between precision and recall?
\begin{enumerate}[label=\Alph*.]
\item Adjust the decision threshold
\item Change the cluster size
\item Change the Delta table version
\item Change the Git branch
\end{enumerate}
\end{frame}
\begin{frame}{Answer 76}
\textbf{Correct: A} --- a higher threshold typically raises precision and lowers recall.\\
\vspace{0.5em}
\textbf{Wrong:}
\begin{itemize}
\item B: unrelated.
\item C: unrelated.
\item D: unrelated.
\end{itemize}
\end{frame}

% ---------------------------------------------------------------
\begin{frame}{Question 77}
A batch scoring job needs to write predictions to Delta daily. Which pattern is typical?
\begin{enumerate}[label=\Alph*.]
\item Call a REST endpoint per row from the driver
\item Email predictions
\item Read input from Delta, apply the model with \texttt{spark\_udf} or \texttt{score\_batch}, then write results to a Delta table
\item Save predictions in notebook variables
\end{enumerate}
\end{frame}
\begin{frame}{Answer 77}
\textbf{Correct: C} --- it is a standard batch inference pipeline.\\
\vspace{0.5em}
\textbf{Wrong:}
\begin{itemize}
\item A: does not scale.
\item B: not a pipeline.
\item D: not durable.
\end{itemize}
\end{frame}

% ---------------------------------------------------------------
\begin{frame}{Question 78}
Why might a team prefer batch inference over real-time serving?
\begin{enumerate}[label=\Alph*.]
\item It does not require a model
\item It always has lower latency
\item It avoids data entirely
\item Predictions are needed periodically and cost and throughput matter more than latency
\end{enumerate}
\end{frame}
\begin{frame}{Answer 78}
\textbf{Correct: D} --- batch is cheaper for large volumes.\\
\vspace{0.5em}
\textbf{Wrong:}
\begin{itemize}
\item A: it does.
\item B: it has higher latency.
\item C: it uses data.
\end{itemize}
\end{frame}

% ---------------------------------------------------------------
\begin{frame}{Question 79}
A real-time use case needs sub-second predictions. Which option fits best?
\begin{enumerate}[label=\Alph*.]
\item A notebook run manually
\item A Model Serving endpoint
\item A Delta \texttt{VACUUM}
\item A nightly job
\end{enumerate}
\end{frame}
\begin{frame}{Answer 79}
\textbf{Correct: B} --- endpoints provide low-latency REST inference.\\
\vspace{0.5em}
\textbf{Wrong:}
\begin{itemize}
\item A: not production-grade.
\item C: unrelated.
\item D: too slow.
\end{itemize}
\end{frame}

% ---------------------------------------------------------------
\begin{frame}{Question 80}
Which workload setting on a serving endpoint affects how many concurrent requests it can handle?
\begin{enumerate}[label=\Alph*.]
\item The notebook theme
\item Workload size (concurrency)
\item The Git branch
\item The MLflow experiment name
\end{enumerate}
\end{frame}
\begin{frame}{Answer 80}
\textbf{Correct: B} --- larger sizes handle more concurrent requests.\\
\vspace{0.5em}
\textbf{Wrong:}
\begin{itemize}
\item A: unrelated.
\item C: unrelated.
\item D: unrelated.
\end{itemize}
\end{frame}

% ---------------------------------------------------------------
\begin{frame}{Question 81}
A team wants to compare a new model against the current one using live traffic safely. What is a good approach?
\begin{enumerate}[label=\Alph*.]
\item Route a small share of traffic to the new version on the same endpoint
\item Replace the production model without testing
\item Test only on the training data
\item Disable monitoring
\end{enumerate}
\end{frame}
\begin{frame}{Answer 81}
\textbf{Correct: A} --- traffic splitting supports gradual rollout.\\
\vspace{0.5em}
\textbf{Wrong:}
\begin{itemize}
\item B: risky.
\item C: not safe.
\item D: removes safety.
\end{itemize}
\end{frame}

% ---------------------------------------------------------------
\begin{frame}{Question 82}
A new model version performs badly in production. What is the fastest rollback when clients load \texttt{@Champion}?
\begin{enumerate}[label=\Alph*.]
\item Delete the new model's experiment
\item Rewrite all clients
\item Retrain from scratch immediately
\item Point the \texttt{Champion} alias back to the previous version
\end{enumerate}
\end{frame}
\begin{frame}{Answer 82}
\textbf{Correct: D} --- clients pick up the older version.\\
\vspace{0.5em}
\textbf{Wrong:}
\begin{itemize}
\item A: does not roll back.
\item B: unnecessary.
\item C: slower.
\end{itemize}
\end{frame}

% ---------------------------------------------------------------
\begin{frame}{Question 83}
What should a CI/CD process for models typically include before promotion?
\begin{enumerate}[label=\Alph*.]
\item Manual copying of files between workspaces only
\item Automated tests and validation against baseline metrics
\item Skipping evaluation
\item Deleting old versions first
\end{enumerate}
\end{frame}
\begin{frame}{Answer 83}
\textbf{Correct: B} --- promotion should be gated by checks.\\
\vspace{0.5em}
\textbf{Wrong:}
\begin{itemize}
\item A: error prone.
\item C: risky.
\item D: prevents rollback.
\end{itemize}
\end{frame}

% ---------------------------------------------------------------
\begin{frame}{Question 84}
Which job setting reruns a task automatically after a transient failure?
\begin{enumerate}[label=\Alph*.]
\item A cluster policy
\item Auto-termination
\item Task retries
\item A model alias
\end{enumerate}
\end{frame}
\begin{frame}{Answer 84}
\textbf{Correct: C} --- retries help with temporary errors.\\
\vspace{0.5em}
\textbf{Wrong:}
\begin{itemize}
\item A: limits configuration.
\item B: a cluster setting.
\item D: labels versions.
\end{itemize}
\end{frame}

% ---------------------------------------------------------------
\begin{frame}{Question 85}
How do you pass a parameter to a notebook task in a job and read it in the notebook?
\begin{enumerate}[label=\Alph*.]
\item Use \texttt{df.cache()}
\item Edit the notebook for each run
\item Define a job parameter and read it with \texttt{dbutils.widgets.get("name")}
\item Use \texttt{dbutils.fs.put}
\end{enumerate}
\end{frame}
\begin{frame}{Answer 85}
\textbf{Correct: C} --- job parameters appear as widgets.\\
\vspace{0.5em}
\textbf{Wrong:}
\begin{itemize}
\item A: not for parameters.
\item B: unnecessary.
\item D: writes files.
\end{itemize}
\end{frame}

% ---------------------------------------------------------------
\begin{frame}{Question 86}
Which tool captures serving request/response data to help monitor model quality?
\begin{enumerate}[label=\Alph*.]
\item Git folders
\item Inference tables
\item Cluster policies
\item \texttt{ZORDER}
\end{enumerate}
\end{frame}
\begin{frame}{Answer 86}
\textbf{Correct: B} --- payloads are logged to a Delta table.\\
\vspace{0.5em}
\textbf{Wrong:}
\begin{itemize}
\item A: version control.
\item C: limit configuration.
\item D: file layout.
\end{itemize}
\end{frame}

% ---------------------------------------------------------------
\begin{frame}{Question 87}
Ground-truth labels arrive weeks after predictions. How can you monitor model quality in the meantime?
\begin{enumerate}[label=\Alph*.]
\item Drop the monitoring tables
\item Retrain every minute
\item Track input drift and prediction distributions, then compute accuracy once labels arrive
\item Do nothing until labels arrive
\end{enumerate}
\end{frame}
\begin{frame}{Answer 87}
\textbf{Correct: C} --- proxy signals detect issues earlier.\\
\vspace{0.5em}
\textbf{Wrong:}
\begin{itemize}
\item A: removes visibility.
\item B: wasteful.
\item D: misses early problems.
\end{itemize}
\end{frame}

% ---------------------------------------------------------------
\begin{frame}{Question 88}
Which change in the input data indicates data drift?
\begin{enumerate}[label=\Alph*.]
\item The experiment name changes
\item The model file size grows
\item The cluster restarts
\item The distribution of a feature in production differs from training
\end{enumerate}
\end{frame}
\begin{frame}{Answer 88}
\textbf{Correct: D} --- inputs have shifted.\\
\vspace{0.5em}
\textbf{Wrong:}
\begin{itemize}
\item A: not drift.
\item B: not drift.
\item C: not drift.
\end{itemize}
\end{frame}

% ---------------------------------------------------------------
\begin{frame}{Question 89}
Why log the training dataset version or Delta table version with each run?
\begin{enumerate}[label=\Alph*.]
\item To reproduce the exact data used to train the model
\item To avoid needing a signature
\item To reduce model size
\item To speed up training
\end{enumerate}
\end{frame}
\begin{frame}{Answer 89}
\textbf{Correct: A} --- reproducibility requires data lineage.\\
\vspace{0.5em}
\textbf{Wrong:}
\begin{itemize}
\item B: not the purpose.
\item C: not the purpose.
\item D: not the purpose.
\end{itemize}
\end{frame}

% ---------------------------------------------------------------
\begin{frame}{Question 90}
Which element makes an MLflow model reproducible on another machine?
\begin{enumerate}[label=\Alph*.]
\item Only the run name
\item The cluster ID
\item Logged dependencies (requirements or conda environment) along with the model artifact
\item Only the notebook theme
\end{enumerate}
\end{frame}
\begin{frame}{Answer 90}
\textbf{Correct: C} --- the environment can be recreated.\\
\vspace{0.5em}
\textbf{Wrong:}
\begin{itemize}
\item A: not enough.
\item B: not portable.
\item D: not related.
\end{itemize}
\end{frame}

% ---------------------------------------------------------------
\begin{frame}{Question 91}
Which Unity Catalog object is the top-level container below the metastore?
\begin{enumerate}[label=\Alph*.]
\item Schema
\item Table
\item Volume
\item Catalog
\end{enumerate}
\end{frame}
\begin{frame}{Answer 91}
\textbf{Correct: D} --- the hierarchy is metastore, catalog, schema, then tables, volumes, models.\\
\vspace{0.5em}
\textbf{Wrong:}
\begin{itemize}
\item A: sits under a catalog.
\item B: sits under a schema.
\item C: sits under a schema.
\end{itemize}
\end{frame}

% ---------------------------------------------------------------
\begin{frame}{Question 92}
What does \texttt{GRANT USE SCHEMA ON SCHEMA main.ml TO group} allow, together with other needed privileges?
\begin{enumerate}[label=\Alph*.]
\item Creating clusters
\item Access to objects within the schema
\item Reading every table automatically
\item Deleting the catalog
\end{enumerate}
\end{frame}
\begin{frame}{Answer 92}
\textbf{Correct: B} --- USE SCHEMA is required to reach objects inside it.\\
\vspace{0.5em}
\textbf{Wrong:}
\begin{itemize}
\item A: not granted.
\item C: SELECT is still needed.
\item D: not granted.
\end{itemize}
\end{frame}

% ---------------------------------------------------------------
\begin{frame}{Question 93}
Who owns a newly created Unity Catalog table by default?
\begin{enumerate}[label=\Alph*.]
\item The cluster
\item The metastore only
\item The user or principal who created it
\item All workspace users
\end{enumerate}
\end{frame}
\begin{frame}{Answer 93}
\textbf{Correct: C} --- owners can manage grants on the object.\\
\vspace{0.5em}
\textbf{Wrong:}
\begin{itemize}
\item A: not an owner.
\item B: not the default owner.
\item D: not by default.
\end{itemize}
\end{frame}

% ---------------------------------------------------------------
\begin{frame}{Question 94}
Why use a service principal to run a production job?
\begin{enumerate}[label=\Alph*.]
\item It removes the need for permissions
\item It avoids tying the job to an individual user's account
\item It replaces Unity Catalog
\item It makes models more accurate
\end{enumerate}
\end{frame}
\begin{frame}{Answer 94}
\textbf{Correct: B} --- service principals suit automation.\\
\vspace{0.5em}
\textbf{Wrong:}
\begin{itemize}
\item A: it still needs permissions.
\item C: it does not.
\item D: not related.
\end{itemize}
\end{frame}

% ---------------------------------------------------------------
\begin{frame}{Question 95}
A column contains sensitive personal data. Which Unity Catalog capability can restrict who sees raw values?
\begin{enumerate}[label=\Alph*.]
\item Column masks or row filters
\item \texttt{mlflow.autolog()}
\item \texttt{OPTIMIZE}
\item \texttt{df.cache()}
\end{enumerate}
\end{frame}
\begin{frame}{Answer 95}
\textbf{Correct: A} --- they apply fine-grained access control.\\
\vspace{0.5em}
\textbf{Wrong:}
\begin{itemize}
\item B: logs runs.
\item C: compacts files.
\item D: does not control access.
\end{itemize}
\end{frame}

% ---------------------------------------------------------------
\begin{frame}{Question 96}
What is the benefit of governing models and features in Unity Catalog?
\begin{enumerate}[label=\Alph*.]
\item Automatic accuracy gains
\item Removal of the need for validation
\item Centralized access control, auditing, and lineage across data and models
\item Free compute
\end{enumerate}
\end{frame}
\begin{frame}{Answer 96}
\textbf{Correct: C} --- one governance layer covers data and ML assets.\\
\vspace{0.5em}
\textbf{Wrong:}
\begin{itemize}
\item A: not provided.
\item B: not provided.
\item D: not provided.
\end{itemize}
\end{frame}

% ---------------------------------------------------------------
\begin{frame}{Question 97}
Which code fills missing values in column \texttt{age} with the median in Spark ML?
\begin{enumerate}[label=\Alph*.]
\item \texttt{Imputer(inputCols=["age"], outputCols=["age\_imp"], strategy="median").fit(df).transform(df)}
\item \texttt{Imputer(strategy="median").transform(df)}
\item \texttt{df.fillna(median)} with no value computed
\item \texttt{StandardScaler(strategy="median")}
\end{enumerate}
\end{frame}
\begin{frame}{Answer 97}
\textbf{Correct: A} --- Imputer is fit and then used to transform.\\
\vspace{0.5em}
\textbf{Wrong:}
\begin{itemize}
\item B: must fit first and set columns.
\item C: no median is computed.
\item D: not a StandardScaler option.
\end{itemize}
\end{frame}

% ---------------------------------------------------------------
\begin{frame}{Question 98}
What input does Spark ML \texttt{StandardScaler} require?
\begin{enumerate}[label=\Alph*.]
\item A pandas Series
\item A vector column, for example from \texttt{VectorAssembler}
\item A Delta table path
\item A string column
\end{enumerate}
\end{frame}
\begin{frame}{Answer 98}
\textbf{Correct: B} --- scalers operate on vector columns.\\
\vspace{0.5em}
\textbf{Wrong:}
\begin{itemize}
\item A: not supported.
\item C: not supported.
\item D: not supported.
\end{itemize}
\end{frame}

% ---------------------------------------------------------------
\begin{frame}{Question 99}
Which call sets the MLflow experiment by workspace path?
\begin{enumerate}[label=\Alph*.]
\item \texttt{mlflow.start\_path(path)}
\item \texttt{mlflow.log\_experiment(path)}
\item \texttt{mlflow.use\_experiment(path)}
\item \texttt{mlflow.set\_experiment("/Users/me@company.com/churn")}
\end{enumerate}
\end{frame}
\begin{frame}{Answer 99}
\textbf{Correct: D} --- it creates or selects the experiment at the path.\\
\vspace{0.5em}
\textbf{Wrong:}
\begin{itemize}
\item A: not an API.
\item B: not an API.
\item C: not an API.
\end{itemize}
\end{frame}

% ---------------------------------------------------------------
\begin{frame}{Question 100}
Which workload most benefits from a GPU cluster?
\begin{enumerate}[label=\Alph*.]
\item Running \texttt{OPTIMIZE} on Delta
\item Training a small scikit-learn linear model
\item Writing a CSV file
\item Training a deep neural network
\end{enumerate}
\end{frame}
\begin{frame}{Answer 100}
\textbf{Correct: D} --- GPUs accelerate large matrix operations.\\
\vspace{0.5em}
\textbf{Wrong:}
\begin{itemize}
\item A: does not need GPUs.
\item B: CPU is enough.
\item C: does not need GPUs.
\end{itemize}
\end{frame}

\end{document}